\section{Introduction}
\label{sec:introduction}

\begin{outline}
  \item Neuroimaging is moving from controlled stimuli to naturalistic ones~\citep{sonkusare2019naturalistic,nastase2020keepitreal}, and from watching to acting: video games give rich, closed-loop tasks in which what the participant sees depends on what they did~\citep{cross2021deep,tomov2023neural,harel2023controller,boyle2025cneuromod}.
  \item A closed-loop stimulus can only be analysed if it can be reconstructed: the run must be replayable, frame by frame, from what was recorded.
  \item Comparing brains with artificial agents needs the same task for both. Reinforcement learning already has a common medium, the environment~\citep{brockman2016gym,towers2024gymnasium}, used for games~\citep{bellemare2013ale,mnih2015dqn} and for neuroscience tasks~\citep{molanomazon2022neurogym}; human experiments are written as scripts for experiment libraries~\citep{peirce2019psychopy,pallier2026goxpyriment}, whose timing has been compared in detail~\citep{bridges2020timing}.
  \item The gap: an environment knows only steps, a script knows only wall time, and neither records the translation between the two; a model's step cannot be aligned with a participant's brain signal unless that translation is part of the data.
  \item Contributions: (1) a framework in which every task is an environment stepped in one of three regimes, with a record that maps engine time to human time (Section~\ref{sec:framework}); (2) software that implements it for 15 game backends, with one record format and bit-exact replay (Sections~\ref{sec:architecture} and~\ref{sec:logging}); (3) a description of the acquisition setup and measurements of its timing that make the record trustworthy (Section~\ref{sec:multimodal}); (4) people and artificial agents on the same environments, with the same seeds and the same record (Section~\ref{sec:players}).
\end{outline}
